\chapter{Recap and Background}
\section{Numerical Relativity}

Numerical Relativity (NR) can be defined as the science of solving the Einstein's equations using numerical methods in such a way that the discretised solutions agree with the continuum solutions at the limit of infinite resolution. They are a set of coupled 4-dimensional Partial Differential Equations (PDEs) that connect the dynamics between matter and the spacetime curvature.

$$
G_{ab}=16 \pi T_{ab}
$$

But its tensorial structure prevents us from approaching them as a usual time evolution PDE problems where we see equations like $\partial_t u =f(x^i,\partial_x u)$. So inorder to cast Einstein's equations into a form suitable for time evolution, we assume 
that the spacetime is globally hyperbolic. This guarantees the existence of a
foliation into space-like Cauchy hypersurfaces $\Sigma_t$, which are the level 
sets of a scalar time function $t : \mathcal{M} \to \mathbb{R}$. Each slice 
admits a unique future-directed timelike unit normal vector \footnote{We use 
abstract index notation following Wald \cite{wald} or Hawking and Ellis~\cite{hawkingellis}.} 
\[
n^a = -\,\alpha \nabla^a t, \qquad 
\alpha^{-2} = -\,g_{ab}\, \nabla^a t\, \nabla^b t,
\]
where $\alpha$ is the lapse function, $g_{ab}$ is the spacetime metric, and 
$\nabla_a$ is the covariant derivative compatible with $g_{ab}$. This 
construction allows the decomposition of Einstein's equations into evolution and 
constraint equations within the $3+1$ formalism. In a coordinate basis adapted 
to the foliation $(\partial_t,\partial_i)$, the vectors $\partial_i$ lie within 
the hypersurfaces $\Sigma_t$, while $\partial_t$ need not be aligned with the 
normal vector $n^a$. Instead, $\partial_t$ generally maps a point on one slice 
to a point with the same coordinate label on the next slice, whereas $n^a$ 
connects orthogonal (normal) directions across slices. The difference is encoded 
in the shift vector $\beta^i$ through the standard decomposition $\partial_t^a = \alpha n^a + \beta^i \partial_i^a,$ which parametrises how spatial coordinates shift from one hypersurface to the 
next.
\begin{figure}
\centering

\begin{minipage}{0.01\textwidth}

\begin{tikzpicture}[scale=0.7]
    \draw[blue,thick,fill=fau_blue_sk] (-1,0) to[out=30, in=200] (5,0)
         to (5.5,1) to [out=200, in=30] (-0.5,1) to (-1,0);
    \node[blue] at (2.25,0.6) {$\Sigma_{t_1}$};
    \node[rose] at (-1.5,0.5) {$t_1$};
    \draw[thick,rose] (-2,0.5) -- (-1.8,0.5);

    \draw[blue,thick,fill=fau_blue_sk] (-1,1.5) to[out=30, in=200] (5,1.5)
         to (5.5,2.5) to [out=200, in=30] (-0.5,2.5) to (-1,1.5);
    \node[blue] at (2.25,2.1) {$\Sigma_{t_2}$};
    \node[rose] at (-1.5,2) {$t_2$};
    \draw[thick,rose] (-2,2) -- (-1.8,2);

    \draw[blue,thick,fill=fau_blue_sk] (-1,3) to[out=30, in=200] (5,3)
         to (5.5,4) to [out=200, in=30] (-0.5,4) to (-1,3);
    \node[blue] at (2.25,3.6) {$\Sigma_{t_3}$};
    \node[rose] at (-1.5,3.5) {$t_3$};
    \draw[thick,rose] (-2,3.5) -- (-1.8,3.5);

    \draw[line width=1.5,rose,->] (-2,-0.5) to (-2,5);
    \node[rose,above] at (-2,5) {Timelike direction};

    \draw[blue,thick,decorate,decoration={brace,mirror,amplitude=10pt}]
                  (6,0) -- (6,4);
    \node[right,blue,align=left] at (6.5,2) {Spacelike \\ hypersurfaces};
\end{tikzpicture}
\end{minipage}
\hfill
\begin{minipage}{0.38\textwidth}
\centering
\begin{tikzpicture}[scale=0.8]

    \draw[line width=1.5,forest_green](1,0.9) ++ (-135:0.7) -- ++ (-135:1.5);
    \draw[line width=1.5,rose] (1,-1) to (1,0.36);

    \draw[blue,thick,fill=fau_blue_sk] (-0.75,0) to[out=30, in=200] (5,0)
         to (5.5,1.1) to [out=200, in=30] (-0.25,1.1) to (-0.75,0);
    \node[blue] at (4.5,0.4) {$\Sigma_{t}$};
    \node[rose] at (-0.75,0.6) {$t$};

    \draw[line width=1.5,rose] (1,0.9) to (1,2.36);
    \node[left,rose] at (1,1.9) {$\alpha n^\mu dt$};

    \draw[line width=1.5,forest_green](1,0.9) ++ (45:0.02) -- ++ (45:1.8);

    \draw[blue,thick,fill=fau_blue_sk] (-0.75,2) to[out=30, in=200] (5,2)
         to (5.5,3.1) to [out=200, in=30] (-0.25,3.1) to (-0.75,2);
    \node[blue] at (4.5,2.4) {$\Sigma_{t + dt}$};
    \node[rose] at (-1.1,2.6) {$t + dt$};

    \draw[line width=1.5,blue,->](1,2.9) 
        to node[above,blue] {$\beta^i dt$} (2.66, 2.56);

    \draw[line width=1.5,rose,dashed] (1,0.4) to (1,0.9);
    \draw[line width=1.5,rose,dashed,->] (1,2.4) to (1,2.9);
    \draw[line width=1.5,rose,->] (1,2.9) to (1,4.5);

    \draw[line width=1.5,forest_green,dashed](1,0.9) ++ (-135:0.02)
                                                 -- ++ (-135:0.7);

    \node[green,above right] at (2.6, 2.2) {$x^i$};
    \node[green,above right] at (1,0.6) {$x^i$};

    \draw[line width=1.5,forest_green,dashed,->]
        (1,0.9) ++ (45:1.75) -- ++ (45:0.65);

    \node[right,forest_green] at ($(1,0.9) + (40:1.15)$) {$t^\mu dt$};

    \draw[line width=1.5,forest_green,->]
        (1,0.9) ++ (45:2.35) 
        to node[above,forest_green] {$t^\mu$} ++ (45:1.75);

    \node[forest_green,above] 
        at ($(1,0.9) + (45:2.35) + (45:1.75)$) {Coordinate line};

    \node[rose,left] at (1,4.2) {$n^\mu$};
    \node[rose,above] at (1,4.5) {Normal line};

\end{tikzpicture}

\end{minipage}

\caption{
(a) Spacelike hypersurfaces sliced along the timelike direction. 
(b) Relation between slices showing lapse $\alpha$ and shift $\beta^i$.
}
\end{figure}

By using the orthogonal projection operator to the hypersurface $\perp_{ab} = \delta_{ab} +n_a n_b$ we can project all  the relevant quantities that we require such as the metric, Riemann curvature tensor etc. This projected part of $g_{ab}$ is called spatial metric $\gamma_{ab}$ and D is the derivative associated with this metric. We define extrinsic curvature $K_{ab}$ to be $-\gamma^c_a \nabla_c n_b$ which in vague terms give the information of the time derivative of the spatial metric by the relation $K_{ab} = -\mathcal{L}_n \gamma_{ab}$ where $\mathcal{L}$ is the lie derivative with respect to a vector field. Using all these we can decompose the Einstein's equations into a set of constraint equations and evolution equations as given below.
\begin{equation}
\label{eq:Constraints}
R + K^2 - K_{ab} K^{ab} = 16\pi \rho 
\qquad
D_b K^{ab} - D^a K = 8\pi j^a,
\end{equation}
which correspond to Hamiltonian and momentum constraints respectively, R corresponds to Ricci scalar associated with the $\gamma_{ab}$, $j^a =- \perp T_{ab}n^b, s_{ab}= \perp T_{ab}, \mathcal{\rho} =T_{ab}n^a n^b $ \footnote{$\perp$ means using projection operator on every free index}

\[
\partial_t \gamma_{ij} 
= -2\alpha K_{ij} + \mathcal{L}_{\beta}\gamma_{ij},
\]
\[
\partial_t K_{ij} 
= - D_i D_j \alpha 
  + \alpha( R_{ij} 
  - 2 K_{ik} K^{k}{}_{j} 
  + K K_{ij})
  + 4\pi \alpha \!\left[\, \gamma_{ij}(s - \rho) - 2 s_{ij} \right]
  + \mathcal{L}_{\beta} K_{ij}.\footnote{where i,j $-> $1,2,3 (Spatial)}
\]
 Now we can see that the GR equations have been formulated into a way which is ideal for time evolution using standard PDE techniques(Since all the elements in the RHS are spatial). However, more modern formulations such as BSSN, Z4C, and related systems include constraint-damping mechanisms that improve numerical stability, since in a discretised scheme the constraints cannot be satisfied exactly. It is also well known that the constraint subsystem is closed, meaning that the time derivatives of the constraints depend only on the constraints themselves. Consequently, if the constraint equations are solved initially, they remain satisfied for a finite time during the evolution. For further details, see \cite{hilditch_2024}.

\section{Critical Phenomena in Gravitational Collapse}

\section*{Recap and Motivation}

\section{Horizons}

The main tool we use to classify spacetimes to subcritical or supercritical is by probing to find horizons, specifically Apparent Horizons (AHs) in the evolved slice. The main advantage of AH is that, unlike event horizons they are not \textit{teleological} i.e we dont need the information about the complete spacetime to get the information of apparent horizon. In numerically computed dynamical spacetimes like ours, we might not be in the position to evolve it completely inorder to classify them, hence AHs are the right tool to understand them.

But what is an AH and why does it relate to the notion of a singular spacetime?

Essentially an AH can be defined as the outermost, Marginally Outer Trapped Surface (MOTS) which are smooth closed 2-dimensional submanifold S in the spacetime slice that satisfies the condition that the expansion of the outgoing null geodesics vanishes at each point in S. The expansion is defined as 

$$
\Theta = h^{ab} \nabla_a \zeta_b \qquad \text{where} \; h^{ab}= g^{ab}-s^as^b + n^a n^b 
$$

$h^{ab}$ is the induced metric of the 2-dimensional submanifold, $s^a$ is its normal vector and $\zeta$ is the tangent vector along the outgoing null geodesic . Solving this we get the equation of the form,

\begin{equation}
\Theta = -K + D_i s^i +s^is^j K_{ij} 
\label{Horizon}
\end{equation}

Hence, our aim is to find surfaces satisfying ($\Theta = 0$), and we identify the apparent horizon (AH) as the \textbf{outermost marginally outer trapped surface (MOTS)}.

The importance of apparent horizons is closely tied to the singularity theorems of general relativity. These theorems state that, under suitable energy conditions (such as the \textit{null energy condition}, which is satisfied in vacuum since ($T_{ab}=0$)) and global hyperbolicity (satisfied since the initial data is constraint satisfying), the presence of a trapped surface implies that the spacetime is geodesically incomplete, i.e., contains a singularity.

Thus, the formation of a MOTS or apparent horizon serves as a strong indicator of black hole formation in the spacetime’s future development. This provides a practical and geometric criterion for classifying spacetimes in our analysis.


\section{Gauge Freedom in GR (NR)}

One of the most important characteristic of GR is its so called diffeomorphism invariance which essentially means
that the "physics" inside theory is unchanged under smooth coordinate transformations. In other words, different coordinate systems (or gauges) can describe the same underlying spacetime geometry.


Mathematically, this invariance implies that if $g_{\mu\nu}(x)$ is a solution to Einstein's equations, then so is the transformed metric under a diffeomorphism $x^\mu \rightarrow x'^\mu(x)$, with the metric transforming as
\begin{equation}
g'_{\mu\nu}(x') = \frac{\partial x^\alpha}{\partial x'^\mu} \frac{\partial x^\beta}{\partial x'^\nu} g_{\alpha\beta}(x).
\end{equation}

Often in theoretical studies we utilise various gauges to uncover the internal dynamics of the spacetime variables. The easiest example would in the case of gravitational waves where we often assume a linearised gravity ($g_{\mu\nu}= \eta_{\mu \nu} + h_{\mu \nu}$ where $O(h) << 1$) where the Minkowskian flat space metric is considered as the background metric of the spacetime.  Here to retrieve the notion of gravitational waves we use something known as Lorenz gauge or De Donder gauge ($\partial^{\mu} \bar{h}_{\mu\nu}$) where $\bar{h}$ is the trace reversed metric perturbation. 
Using this gauge we will be able to see that, the metric perturbation has formalised into the form of a wave equation.

$$
\Box \bar{h}^{\mu\nu}= -16 \pi T^{\mu \nu}
$$
This is finest and easiest example of utilising gauges to understand the internal physics. But in the case of the NR, it is utilised differently. As
we know in NR, we try to see Einstein's equations as a set of PDE's to solve numerically. But to solve a problem numerically it is essential that the problem we are solving is \textit{well-posed} and devoid of any coordinate artifacts that coordinate system may evolve into(coordinate singularities , constraint violation and so on). So essentially, we utilise the gauge freedom to counter these problems. Based on how we tackle them has given rise to various formulations of NR like BSSN, Z4C etc. In our case we use the Generalised Harmonic Formulation of NR.

\section{Generalised Harmonic Gauge}

The main advantage of the Generalised Harmonic gauge is that it is, by construction, a well-posed formulation of Einstein's equations. This makes it particularly suitable for numerical evolution, as it provides a strongly hyperbolic system with well-defined characteristic structure. Due to this property, practical considerations aside, in theory they are ideal for any kind of numerical evolutions performed with GR in mind.  The well-posedness of this formulation comes from the trace reversed form of Einstein's equation i.e



\begin{equation}
R_{\mu\nu}
= -\frac{1}{2} g^{\alpha\beta} \partial_\alpha \partial_\beta g_{\mu\nu}
+ \partial_{(\mu} \Gamma_{\nu)}
+ \Gamma^\alpha \Gamma_{\mu\nu\alpha}
- \Gamma^\alpha_{\mu\beta} \Gamma^\beta_{\nu\alpha}
\label{GHG}
\end{equation}

where $
\Gamma^\mu = g^{\alpha\beta} \Gamma^\mu_{\alpha\beta}$ called the contracted Christophel symbols. The well-posedness or ill-posedness of a PDE is often determined by its principal terms (Terms containing the highest derivatives). In our problem, the principal terms are the first and the second part of the Equation \ref{GHG} since they contain the second derivatives of the metric. But since the wave equation is a well posed problem, the problematic part of the equation is the second term (derivative of contracted Christophels). But for that we can utilise the gauge freedom to ensure that they don't contain second derivatives of the metric i.e
$\Gamma^a=\nabla_b \nabla^b x^a=-H^a(x,g)$ where $H^a$ is called the Gauge Source function and it is a function of atmost the metric and not its derivatives.

This gives us the following equations of motion for the lapse($\alpha$) and shift($\beta$)

\begin{equation}
\partial_t \alpha = -\alpha^2 (K + H_n) + \mathcal{L}_\beta \alpha
\qquad \partial_t \beta^i = \alpha^2 (\Gamma^i + \perp H^i) - \alpha D^i \alpha + \beta^j \partial_j \beta^i
\label{lapse=shift}
\end{equation}
where $H^n= n^a H_a$

The $H^a$'s are freely specifiable functions but as we mentioned before, we choose it in a way to have proper evolutions without any coordinate pathologies. That choice will be mentioned in the coming sections.
% Remember to input this to the presentative tex file before compiling.